\documentclass[xelatex,ja=standard,a4paper,10pt]{bxjsarticle}
\geometry{margin=18mm,top=20mm,bottom=22mm,columnsep=7mm}
\usepackage{amsmath,amssymb,mathtools,bm}
\usepackage{graphicx}
\usepackage{booktabs}
\usepackage{multirow}
\usepackage{tabularx}
\usepackage{array}
\usepackage{siunitx}
\usepackage{xcolor}
\usepackage{tikz}
\usetikzlibrary{arrows.meta,positioning,fit,shapes.geometric}
\usepackage{enumitem}
\usepackage{caption}
\usepackage{subcaption}
\usepackage{float}
\usepackage{placeins}
\usepackage{url}
\usepackage[
  colorlinks=true,
  linkcolor=blue!45!black,
  citecolor=blue!45!black,
  urlcolor=blue!55!black
]{hyperref}
\ifpTeX
  \usepackage{pxjahyper}
\fi
\usepackage[nameinlink,noabbrev]{cleveref}

\captionsetup{font=small,labelfont=bf}
\setlength{\emergencystretch}{2em}
\setlength{\parindent}{1em}
\setlength{\parskip}{0pt}
\setlist{nosep,leftmargin=1.8em}
\sisetup{
  detect-all,
  group-separator={,},
  group-minimum-digits=4,
  range-phrase={--},
  range-units=single
}

\newcommand{\pt}{p_{\mathrm{T}}}
\newcommand{\met}{E_{\mathrm{T}}^{\mathrm{miss}}}
\newcommand{\Htt}{H\rightarrow\tau^{+}\tau^{-}}
\newcommand{\Ztt}{Z\rightarrow\tau^{+}\tau^{-}}
\newcommand{\epssig}{\varepsilon_{\mathrm{sig}}}
\newcommand{\epsbkg}{\varepsilon_{\mathrm{bkg}}}
\DeclareSIUnit{\GeV}{GeV}
\DeclareSIUnit{\TeV}{TeV}
\newcommand{\code}[1]{\texttt{\detokenize{#1}}}
\newcommand{\todo}[1]{\textcolor{red!70!black}{[要確認: #1]}}

\crefname{figure}{図}{図}
\crefname{table}{表}{表}
\crefname{section}{節}{節}
\crefname{equation}{式}{式}

\hypersetup{
  pdftitle={再構成されたタウ対による等質量H/Z分離の探索},
  pdfauthor={Ryunosuke Baba}
}


\title{等質量 \(H/Z\rightarrow\tau\tau\) 分離解析\\
\large Analysis Note：根拠・仮定・レビュー記録}
\author{Ryunosuke Baba}
\date{2026年7月30日}

\begin{document}
\maketitle

\section{目的と判定状態}

本解析の目的は、質量を約 \SI{91.2}{\GeV} に揃えた
\(H\rightarrow\tau^+\tau^-\) と \(Z\rightarrow\tau^+\tau^-\) の
専用Monte Carlo sampleについて、再構成されたevent、\(\tau\)、track、PFOを
用いたTransformerの分離能を調べることである。
検証する作業仮説は、truth親ボソン\(\pt\)をclass間でmatchingした後にも、
再構成レベルの多変量入力にH/Z labelを識別する情報が残る、というものである。
ここでmatchingとは各分割・\SI{20}{\GeV} bin内のH/Z事象数を
等しくする処理であり、bin内shapeまで同一にする操作ではない。

\begin{description}
  \item[レポート作成状態] 実行済み結果を記述する準備が整っている。
  \item[物理結論の状態] \textbf{Blocked（spin-only解釈）}。
    spin on/off比較、feature ablation、複数seed、systematic variationが
    未実施であるため、観測したAUCをスピン相関だけに帰属させない。
  \item[許容する結論] 固定partial snapshotにおける、
    再構成レベルH/Z分離の探索的proof of concept。
\end{description}

\section{サンプル、selection、評価設計}

入力は \(\sqrt{s}=\SI{13}{\TeV}\) の専用H/Z sampleであり、truth質量は
それぞれ \SI{91.1872}{\GeV} と \SI{91.1876}{\GeV} である。
固定snapshotは生成予定量の63.1\%に相当する。
selection効率の分母はntuple生成時に実際に読み取ったraw event数、
分子はselection後のentry数とする。

eventは決定的hashでtrain/validation/testへ80:10:10に分割してから、
split内かつtruth親ボソン\(\pt\)の\SI{20}{\GeV} bin内でH/Zを同数に
downsampleした。HPO、epoch選択、最終checkpointの固定にはvalidationだけを使い、
最終checkpointをtestへ適用した。ただし同じtest splitは先行モデルにも
使用済みであるため、解析全体から独立したblind holdoutとは扱わない。
test上の信号効率70\%点はtest ROCそのものから記述的に選んだ値であり、
独立に固定したworking pointとは扱わない。

\section{Evidence Ledger}

\small
\noindent
\begin{tabularx}{\textwidth}{@{}p{0.18\textwidth}X
  >{\raggedright\arraybackslash}p{0.22\textwidth}@{}}
  \toprule
  種別 & 本文へ用いる根拠 & 主なsource \\
  \midrule
  observed &
  Hは331 DAOD、661,964 raw、332,932 selected、
  Zは300 DAOD、599,955 raw、306,731 selected。 &
  \path{MakeNtuple/outputs/fixed-partial-v2/summary.json} \\
  observed &
  全64 jobがsuccess、schemaは84 branch、truth親PDG不一致は0。 &
  fixed snapshot summaryとjob metadata \\
  decision &
  selectionは\(\pt\geq\SI{20}{\GeV}\)、\(|\eta|\leq2.5\)、
  \(|q|=1\)、1/3 track、RNN Loose、exactly two、opposite sign。
  truth matchとtrigger requirementは置かない。 &
  \path{MakeNtuple/README.md}、selection実装 \\
  observed &
  matching後はtrain 456,720、validation 57,038、test 56,338で、
  各split内のH/Zは同数、合計570,096 event。 &
  matching manifest \\
  calculated &
  truth親\(\pt\)だけを使うclassifierのAUCはmatching前0.570524、
  matching後0.500660。 &
  matching report \\
  decision &
  final modelはtrial 5の設定で40 epochを最初から再学習し、
  validation 3-epoch平均が最大の23--25の中央、epoch 24を固定。 &
  final selection metadata、history \\
  observed &
  epoch 24のvalidation AUCは0.604712、lossは0.677076。
  test AUCは0.608814、BCE lossは0.674915。 &
  final metrics JSON \\
  calculated &
  test ROCで信号効率0.700025となる記述点は、
  Z効率0.548972、Z除去率1.821586、threshold 0.465762。 &
  final test metrics JSON \\
  calculated &
  testのtruth親\(\pt\) bin別AUC誤差棒はclass-stratified
  bootstrap 1,000回の68\% interval。global intervalではない。 &
  \path{pt_binned_auc_test.pdf}の生成metadata \\
  observed &
  最終train splitのZ event 1件に
  track \(\pt=\SI{7.74e10}{\GeV}\) の非物理的外れ値が残る。 &
  processed shard監査とsource entry照合 \\
  \bottomrule
\end{tabularx}
\normalsize

\section{仮定と既知の制約}

\begin{enumerate}
  \item selectionは再構成\(\tau\)候補へ適用しており、truth-matched候補数、
    真のhadronic-\(\tau\)対、fake候補のH/Z別組成を監査していない。
  \item sampleはpartialかつnominal-onlyであり、断面積、積分luminosity、
    event weight、scale factor、systematic variationを用いない。
  \item matchingしたのはtruth親\(\pt\)だけであり、production radiation、
    20 GeV bin内shape、非単調な\(\pt\)依存、\(\eta\)、recoil、
    detector response、decay-mode組成の差は個別に除いていない。
  \item 最終学習はseed 42の1回で、AUCのglobal statistical intervalもない。
  \item 同じtest splitはHPO直後のモデルと40-epoch最終モデルにそれぞれ1回
    適用済みである。ただし最終checkpointはvalidation-onlyで固定し、
    testを見た後の再選択はしない。
  \item MetMaker error message、invalid tau-track link warning、
    上記の非物理的track外れ値は、最終性能への影響を未評価である。
  \item split keyはROOT basenameとfile-local entryへ依存するため、
    fileのrename/rechunkに不変ではない。
\end{enumerate}

\section{独立レビューと反映記録}

数値、物理解釈、論文形式を独立にレビューした。
event数、split/matching数、HPO、epoch 24、AUC/loss、効率、
主要hashについて、本文と固定成果物の数値不一致は見つからなかった。
一方、物理レビューはspin-only解釈を\textbf{Blocked}、論文形式レビューは
\textbf{Major revision}と判定した。

本文だけで修正可能な指摘には、次を反映した。
\begin{itemize}
  \item ``test''を再利用済みの固定評価分割と定義し、blind性能とは呼ばない。
  \item \(\pt\)整合の保証を20 GeV binごとの事象数一致に限定し、
    raw-\(\pt\) AUCの診断能力にも限界を明記する。
  \item selection効率をraw DAOD事象からの処理通過率と定義し、
    truth \(\tau\)組成の未監査を主要制約へ置く。
  \item HPOの設計と観測結果をMethods/Resultsへ分離し、目標100 trial、
    時刻停止後21 trial、停止閾値、事後的な順位規則を明記する。
  \item sigmoid出力を確率ではなくscoreと呼び、inclusive統計区間がないため
    validation/testのcompatibilityを主張しない。
  \item 全入力特徴量、model profile、software、hashを付録へ置き、
    生成段階の設定不足から再現性の範囲を固定DAOD以降に限定する。
\end{itemize}

文章修正では解消できず、追加解析を要する項目は、truth-matched
hadronic-\(\tau\)組成とfake組成の監査、より厳密なkinematic control、
未使用holdout、inclusive metricの区間、複数seed、spin on/off、
feature ablation、systematic variation、MET/link/outlierの影響評価である。
したがって、論文として成立する範囲は固定partial Monte Carlo標本に対する
再構成候補分類の探索的概念実証までであり、spin感度測定または公開論文相当の
最終物理結果には到達していない。

\section{実装・検証結果}

日本語二段組の本文、構成・図表計画、レビュー前草稿、独立レビュー記録を
TeXで作成した。本文ではMethods/Resultsの分離、二段の解析フロー図、
自立した図表caption、再現性付録、一次文献を追加した。
XeLaTeXで全TeX文書を再buildし、Ghostscriptで本文全ページを画像化して、
重なり、切れ、図表の順序を確認した。remote repositoryでは
\path{/report/}が\path{.gitignore}により無視され、
作業前後でHEADが不変であること、すなわちcommitを作成していないことを確認した。

\end{document}
